\chapter{Grade 10: Analysis of Decisions}\label{sec:grade10-term2-overview}

\section{Learning Strand}\label{sec:grade10-term2-strand}

Complex and non-standard problems solving.

\section{Learning Objective}\label{sec:grade10-term2-objective}

Evaluate all possibilities in a decision-making problem using different strategies.

\section{Assessment Criteria}\label{sec:grade10-term2-criteria}

\begin{enumerate}
	\item Describes the way a problem can be formulated for optimal decision making.
	\item Interprets the parts of a problem such as decision alternatives, fortuitous events, consequences, and states of nature.
	\item Builds a payoff table from a description of the problem.
	\item Explains how a decision-making problem without probabilities works (criterion maximax).
	\item Summarises how a decision-making problem without probabilities works (criterion maximin).
	\item Develops decision-making strategies to address problems with probabilities and the maximum opportunity criterion.
	\item Uses probabilities and expected value to work on a decision-making problem.
	\item Compares the Bayes decision rule for decision making without experimentation and the Bayes, maximin and minimax methods.
	\item Establishes the tree diagram for a decision-making problem.
	\item Evaluates all possibilities in a decision-making problem using the different strategies presented so far.
\end{enumerate}

\section{Key Topics}\label{sec:grade10-term2-topics}

\begin{itemize}
	\item Random variable.
	\item Probability distributions.
	\item Subjective probability.
	\item Mean and variance of discrete distributions.
	\item Investment portfolio.
	\item Expected value.
\end{itemize}

\vspace{6pt}
\noindent\hyperref[table-of-contents]{Back to Top}
